\documentclass[11pt]{article}
\usepackage[margin=1in]{geometry}
\usepackage{amsmath,amssymb,mathtools}
\usepackage{fontspec,unicode-math}
\setmainfont{Latin Modern Roman}
\setsansfont{Latin Modern Sans}
\setmathfont{Latin Modern Math}
\usepackage{xurl,hyperref}
\hypersetup{colorlinks=true,urlcolor=blue,linkcolor=blue}
\newcommand{\sref}[2]{\hyperlink{source:#1:#2}{[S#2]}}
\newcommand{\cataloguescope}[1]{\paragraph{Recorded mathematical question.}#1}
\begin{document}
% UnsolvedMath ID 10000064; source catalogue code AMR-099-0064
\setcounter{section}{10000063}
\section{Accumulation points of packings of $\mathbb{Z}^3$}\label{10000064}
{\small\sffamily UnsolvedMath ID 10000064 · Geometry}
\subsection{Definitions and mathematical statement}

\paragraph{Shared notation.}$\mathbb N=\{1,2,\ldots\}$, and $\mathbb Z,\mathbb Q,\mathbb R,\mathbb C$ denote the integers, rationals, reals and complex numbers. Any other conventions are specified locally.


A sphere packing of the nearest-neighbor lattice graph $\mathbb Z^3$ consists of balls $B(c_v,r_v)\subset\mathbb R^3$, $r_v>0$, with disjoint interiors, such that two balls are tangent exactly when their indices differ by a standard unit vector. A point of the one-point compactification $\mathbb R^3\cup\{\infty\}$ is an accumulation point of the packing if each neighborhood of it meets infinitely many distinct balls.
\paragraph{Statement recorded in the supplied source.}
Prove that every such packing has at most one accumulation point in $\mathbb R^3\cup\{\infty\}$. \sref{10000064}{2}
\subsection{Short English statement}
Can a sphere packing realizing the three-dimensional lattice accumulate anywhere other than a single point, including infinity?
\subsection{Sources}
\begin{description}
\item[\hypertarget{source:10000064:1}{S1}] ULAM.AI and the UnsolvedMath Contributors, \emph{UnsolvedMath: A Curated Collection of Open Mathematics Problems}, record 10000064 (AMR-099-0064). CC BY 4.0. \url{https://huggingface.co/datasets/ulamai/UnsolvedMath}.
\item[\hypertarget{source:10000064:2}{S2}] Itai Benjamini, \emph{Euclidean vs. Graph Metric}, version dated 16 July 2012; original author manuscript preserved in Arquivo.pt capture of 31 December 2020, Section 6, Question 6.5, p. 8. \url{https://arquivo.pt/noFrame/replay/20201231041538id_/http://www.wisdom.weizmann.ac.il/~itai/erd100.pdf}.
\end{description}
\label{end:10000064}
\end{document}
